% This ensures VS Code compiles the main file even when you're editing a chapter.
% !TeX root = ../main.tex

\chapter{Theoretical Background}
\label{cap:back}

This chapter presents the theoretical concepts that support the development of the proposed
system. The system consists of a drone payload, or pod. The test prototype will consist of
two system and a communication channel: The drone payload and a drone system, that will serve
to test the payload's performance. The drone itself will consist of a small embedded computer
that will using a computer vision subsystem to identify a ground target. Then, the drone
system will send the direction angles to the drone payload for it to execute the measurement
and error correction routines. The resulting values are then sent back to the drone system
for logging.

The pod will contain a light-based time-of-flight distance sensor mounted on a gimbal mechanism.
When the drone system sends a specific direction, the pod will measure the distance from the
drone to the target taking  geometric effects such as line-of-sight and parallax into account.

\section{Drone subsystem}

The drone system is responsible for perception, processing, and identification of the initial
payload orientation. In this work, the drone platform incorporates an onboard computing system
capable of running computer vision algorithms. The purpose of this subsystem is to identify
objects of interest in the environment and determine the appropriate pointing direction for
the distance measurement sensor.

The information extracted from image processing algorithms can be used to generate control
commands for the gimbal system. These commands orient the distance sensor toward the
detected target, enabling the measurement of the distance between the drone and the object
of interest.

The drone subsystem therefore performs three primary functions:

\begin{itemize}
  \item Image acquisition from the onboard camera.
  \item Image processing and feature detection.
  \item Target localization and gimbal pointing control.
\end{itemize}

\subsection{Computer vision}

Computer vision is a field of study that enables machines to interpret and understand visual
information from the environment. It involves acquiring images, processing them, extracting
meaningful features, and making decisions based on those features. A typical computer vision
pipeline consists of several stages \cite{cv1}:

\begin{enumerate}
  \item Image acquisition
  \item Pre-processing
  \item Feature extraction
  \item Object detection or classification
  \item Decision or control output
\end{enumerate}

\subsubsection{Image acquisition}

Image acquisition consists of capturing an image from an image sensor, like a camera. cameras
have many different parameters, but the basic ones are \cite{photo}:

\begin{itemize}
  \item \gls{fov}
  \item Resolution
  \item \gls{fps}
\end{itemize}

\gls{fov} defines the angular area the camera can capture. A wider \gls{fov} covers more
of the scene, while a narrower \gls{fov} focuses on a smaller area with greater detail.
Resolution refers to the number of pixels in an image. Higher resolution provides more
detail but increases data size and processing requirements. \gls{fps}, or frame rate,
indicates how many images are captured per second. Higher frame rates produce smoother
motion, while lower frame rates reduce data and processing demands, a simple example can
be seen in Figure~\ref{fig:2img1}.

\begin{figure}[ht]
  \centering
  \includegraphics[width=0.5\textwidth]{2/img1}
  \caption{Field of View of a camera and its resolution}
  \label{fig:2img1}
\end{figure}

\subsubsection{Pre-processing and feature extraction}

Image pre-processing may involve operations such as filtering, noise reduction, and contrast
enhancement. Feature extraction methods detect relevant structures in the image, such as
edges, corners, shapes, or textures. This step usually involves a mask, that is used to
identify objects or track their movement within the camera frame \cite{cv2}.

\subsubsection{Object detection and decision making}

Once a target is detected, its position in the image coordinates can be calculated. This
position can then be transformed into a direction relative to the camera frame. For this,
the inclination of the camera and camera lens angles are taken into account. The direction
information can be sent to the pod, that will point the gimbal-mounted sensor, allowing
the sensor to align with the target detected by the vision system \cite{cv3}.

\subsubsection{OpenCV}

\gls{opencv} (logo in Figure~\ref{fig:2opencv}) is an open-source software library widely
used for computer vision and image processing applications\cite{opencv}. It provides a
comprehensive set of algorithms optimized for real-time performance. The library supports
a wide range of functionalities, including:

\begin{itemize}
  \item Image filtering and transformations
  \item Feature detection and extraction
  \item Object detection and recognition
  \item Camera calibration
  \item Motion analysis
  \item Machine learning integration
\end{itemize}

\begin{figure}[ht]
  \centering
  
\includegraphics[width=0.2\textwidth]{2/opencv}
  \caption{OpenCV's Logo}
  \label{fig:2opencv}
\end{figure}

\gls{opencv} is designed to be highly efficient and supports multiple programming languages,
including C++, Python, and Java. Its optimized implementations make it suitable for embedded
systems and single-board computers.

In this project, \gls{opencv} can be used to process images captured by the drone's camera
in real time. The algorithms implemented in \gls{opencv} can detect objects or features of
interest and determine their location in the image frame. This information is then used to
compute the angular offset between the camera's optical axis and the detected target.

\section{Distance measuring pod}

The distance measuring pod is a payload mounted on the drone that will contain the hardware
responsible for measuring the distance to a target. The pod integrates two main components:
a \gls{tof} distance sensor and a gimbal mechanism that controls the orientation of the sensor.

The purpose of the pod is to measure the distance between the drone and an object detected
by the vision system. Since the camera and the distance sensor may not always share the
same orientation, a gimbal mechanism is used to align the sensor with the detected target.

The pod therefore acts as an independent measurement subsystem that receives pointing
commands from the drone's main processing unit and will not only measure the distance
but also compensate for any geometric error.

\subsection{Light-based time-of-flight sensors}

\gls{tof} distance measurement is a technique used to determine the distance between a sensor
and an object by measuring the time required for light to travel to the object and return to
the sensor \cite{tof1}.

In a typical light-based ToF system, the sensor emits a short pulse of light, usually in
the infrared spectrum. This light propagates through space until it reaches an object,
where it is reflected back toward the sensor. The sensor then detects the reflected light
and measures the time interval between emission and detection. Because the speed
of light is known and constant in air, the distance to the object can be calculated using
the relation~\ref{eq:distance}:

\begin{equation}
  \label{eq:distance}
  d = \frac{c \cdot t}{2}
\end{equation}

%where:

\begin{itemize}
  \item $d$ is the distance to the object
  \item $c$ is the speed of light
  \item $t$ is the measured round-trip travel time
\end{itemize}

The division by two is necessary because the measured time corresponds to the round-trip
travel of the light signal.

\gls{tof} sensors offer several advantages compared to other distance measurement
techniques \cite{tof2}. They can operate independently of ambient lighting conditions, provide
relatively fast measurements, and offer good accuracy over short to medium distances.
However, several factors may affect measurement accuracy , including:

\begin{itemize}
  \item Surface reflectivity of the target
  \item Ambient infrared interference
  \item Sensor field of view
  \item Multipath reflections
\end{itemize}

\subsection{Sensor gimbal}

A gimbal is a mechanical structure that allows an object to rotate about one or more axes
while maintaining controlled orientation. An example of a gimbal stabilizing a camera can be
seen in Figure~\ref{fig:2gimbal}. Gimbal mechanisms are commonly used in camera
stabilization systems, aerospace instruments, and robotic platforms \cite{gimbal}.

\begin{figure}[ht]
  \centering
  \includegraphics[width=0.5\textwidth]{2/gimbal}
  \caption{Stabilizing gimbal with camera}
  \label{fig:2gimbal}
\end{figure}

In the proposed system, the gimbal will allow the distance sensor to be oriented independently
of the drone's body orientation. This capability is important because the drone may need to
maintain a stable flight attitude while directing the sensor toward different targets.

A typical gimbal used in robotic applications provides two or three degrees of freedom. A
two-axis gimbal allows rotation in pitch and yaw, enabling the sensor to point in any
direction within a hemispherical field of view. A three-degrees gimbal include roll control
for full stabilization and orientation control. Ideally, when mounting a sensor on a gimbal,
it is desirable to mount the sensor in the gimbal in a way that its position does not vary
with any change in its angle. This ensures the sensor's position remains constant throughout
all gimbal movements, simplifying parallax compensation \cite{robotics2}.

The gimbal is usually actuated by servo motors, using either brushed or brushless motors.
The control system receives commands from the drone's processing unit and adjusts the
orientation of the sensor accordingly \cite{servos}.

The relationship between the target position detected by the camera and the gimbal angles
must be carefully calibrated. This calibration ensures that the direction indicated by
the computer vision system corresponds to the physical direction in which the sensor must
point.

\section{Line-of-sight and parallax}

Accurate distance measurement requires a clear geometric understanding of the relationship
between the camera, the distance sensor, and the observed target. Two important concepts
in this context are \gls{los} and parallax.

\gls{los} refers to the direct path between the sensor and the target. A valid
measurement can only be obtained if the emitted signal travels unobstructed to the target
and the reflected signal returns to the sensor. The first \gls{los} that must be accounted
for is the one between the camera and the target. The second \gls{los} is the one between
the \gls{tof} sensor and the target \cite{parallax1}.

In systems where the camera and the distance sensor are located at different positions,
their respective lines of sight may not coincide. This spatial separation introduces a
geometric effect known as parallax. An example where both camera and distance sensor are
pointing at the same angle but not at the same target can be seen in Figure~\ref{fig:2img2}.

\begin{figure}[h]
  \centering
  \includegraphics[width=0.5\textwidth]{2/img2}
  \caption{Parallax error happens because distance sensor and camera don't share same LoS}
  \label{fig:2img2}
\end{figure}

Parallax occurs when an object appears at different positions relative to the background
when observed from different viewpoints. In the context of this system, the camera and
the distance sensor may observe the same object from slightly different positions. As a
result, the target detected by the camera may not lie exactly along the measurement axis
of the distance sensor \cite{parallax2}.

The magnitude of the parallax error depends on several factors:

\begin{itemize}
  \item The physical separation between the camera and the distance sensor
  \item The distance to the target
  \item The orientation of the sensor relative to the camera
\end{itemize}

For distant targets, the parallax effect becomes smaller because the angular difference
between the viewpoints decreases. However, for nearby targets the error can become significant.
To reduce parallax effects, the system may use calibration procedures that determine the
spatial transformation between the camera frame and the sensor frame. By knowing this
transformation, it becomes possible to compute the correct pointing direction for the
gimbal-mounted sensor.
